% GENERATED FILE. Edit canonical lessons, then run npm run generate:books.
% canonical-source-hash: fnv1a64:6f5819e84a726e01
% canonical-lessons: GU-C211-hodi, GU-C211-savari-karvi, GU-C211-hankarvu, GU-C211-yatra, GU-C211-vahan

\chapter{Boats, pilgrims and rides}
\label{ch:gu-boats-pilgrims-and-rides}
\hlchaptermodality{211}

\begin{chapteropening}
\textbf{By the end of this chapter:} \emph{I can say how we travelled: by boat, on a ride or on a pilgrimage.}

Complete the last lesson of chapter 211: I can say how we travelled: by boat, on a ride or on a pilgrimage.
\end{chapteropening}

\section[\texorpdfstring{\gu{હોડી} (hoḍī)}{હોડી (hoḍī)}]{\gu{હોડી} (hoḍī) — boat, rowing boat}

\label{lesson:GU-C211-hodi}



\begin{quote}
Before the new one: say the Gujarati for to load, then the Gujarati for to lift.
\end{quote}

\subsection*{You'll want to know: \gu{હોડી}}
\textbf{\gu{હોડી}} — \emph{hoḍī} — \textquotedblleft{}boat, rowing boat\textquotedblright{}.

\begin{quote}
\textbf{\gu{અમે} \gu{હોડીમાં} \gu{નદી} \gu{ઓળંગી}.} — \emph{ame hoḍīmā̃ nadī oḷangī} — \textquotedblleft{}We crossed the river in a boat.\textquotedblright{}
\end{quote}

A feminine noun.

\begin{cousinweb}
\textbf{\gu{હોડી}} (\emph{hoḍī}) is listed by CDIAL, glossed \textquotedblleft{}boat\textquotedblright{}, under the lexical Sanskrit \emph{hōḍa}, \textquotedblleft{}raft, boat\textquotedblright{}, which it takes from Dravidian, as in Kannada \emph{ōḍa}, \textquotedblleft{}boat\textquotedblright{} (CDIAL 14174). DEDR also judges the Indo-Aryan words probably Dravidian (DEDR 1039).
\end{cousinweb}

\begin{grammarlens}[title={What you've built}]
One more word for everyday life at A2.
\end{grammarlens}

\subsection*{Your turn}
\begin{itemize}
\raggedright
  \item \emph{Say it:} \emph{hoḍī}
  \item \emph{Say it:} \emph{hoḍī}, once more
  \item \emph{Say it:} \emph{ame hoḍīmā̃ nadī oḷangī} — We crossed the river in a boat.
  \item \emph{Recall:} say the Gujarati for to load, then the Gujarati for to lift, then say \emph{hoḍī} again
\end{itemize}

\begin{tcolorbox}[breakable,colback=teal!4,colframe=teal!35!black,title={Before you move on}]
What does \gu{હોડી} mean? (\textquotedblleft{}Boat, rowing boat\textquotedblright{}.) Say it once more.
\end{tcolorbox}
\section[\texorpdfstring{\gu{સવારી} \gu{કરવી} (savārī karvī)}{સવારી કરવી (savārī karvī)}]{\gu{સવારી} \gu{કરવી} (savārī karvī) — to ride}

\label{lesson:GU-C211-savari-karvi}



\begin{quote}
Before the new one: say the Gujarati for to lift, then the Gujarati for boat.
\end{quote}

\subsection*{You'll want to know: \gu{સવારી} \gu{કરવી}}
\textbf{\gu{સવારી} \gu{કરવી}} — \emph{savārī karvī} — \textquotedblleft{}to ride\textquotedblright{}.

\begin{quote}
\textbf{\gu{બાળકોએ} \gu{હાથીની} \gu{સવારી} \gu{કરી}.} — \emph{bāḷkoe hāthīnī savārī karī} — \textquotedblleft{}The children rode an elephant.\textquotedblright{}
\end{quote}

\begin{cousinweb}
\textbf{\gu{સવારી} \gu{કરવી}} (\emph{savārī karvī}) is built from \emph{savārī}, \textquotedblleft{}a ride\textquotedblright{}, a loan usually traced to Persian \emph{savārī}, from \emph{savār}, \textquotedblleft{}rider\textquotedblright{}, and the light verb \emph{karvũ}, \textquotedblleft{}to do\textquotedblright{}.
\end{cousinweb}

\begin{grammarlens}[title={What you've built}]
One more word for everyday life at A2.
\end{grammarlens}

\subsection*{Your turn}
\begin{itemize}
\raggedright
  \item \emph{Say it:} \emph{savārī karvī}
  \item \emph{Say it:} \emph{savārī karvī}, once more
  \item \emph{Say it:} \emph{bāḷkoe hāthīnī savārī karī} — The children rode an elephant.
  \item \emph{Recall:} say the Gujarati for to lift, then the Gujarati for boat, then say \emph{savārī karvī} again
\end{itemize}

\begin{tcolorbox}[breakable,colback=teal!4,colframe=teal!35!black,title={Before you move on}]
What does \gu{સવારી} \gu{કરવી} mean? (\textquotedblleft{}To ride\textquotedblright{}.) Say it once more.
\end{tcolorbox}
\section[\texorpdfstring{\gu{હંકારવું} (hankārvũ)}{હંકારવું (hankārvũ)}]{\gu{હંકારવું} (hankārvũ) — to steer, to drive}

\label{lesson:GU-C211-hankarvu}



\begin{quote}
Before the new one: say the Gujarati for boat, then the Gujarati for to ride.
\end{quote}

\subsection*{You'll want to know: \gu{હંકારવું}}
\textbf{\gu{હંકારવું}} — \emph{hankārvũ} — \textquotedblleft{}to steer, to drive\textquotedblright{}.

\begin{quote}
\textbf{\gu{તેણે} \gu{હોડી} \gu{હંકારી}.} — \emph{teṇe hoḍī hankārī} — \textquotedblleft{}He steered the boat.\textquotedblright{}
\end{quote}

\begin{cousinweb}
\textbf{\gu{હંકારવું}} (\emph{hankārvũ}) is a tadbhava. CDIAL lists \emph{hakārvũ}, \textquotedblleft{}to drive\textquotedblright{}, under \emph{*hakkārayati}, \textquotedblleft{}calls\textquotedblright{}, a reconstructed form (CDIAL 13941).
\end{cousinweb}

\begin{grammarlens}[title={What you've built}]
One more word for everyday life at A2.
\end{grammarlens}

\subsection*{Your turn}
\begin{itemize}
\raggedright
  \item \emph{Say it:} \emph{hankārvũ}
  \item \emph{Say it:} \emph{hankārvũ}, once more
  \item \emph{Say it:} \emph{teṇe hoḍī hankārī} — He steered the boat.
  \item \emph{Recall:} say the Gujarati for boat, then the Gujarati for to ride, then say \emph{hankārvũ} again
\end{itemize}

\begin{tcolorbox}[breakable,colback=teal!4,colframe=teal!35!black,title={Before you move on}]
What does \gu{હંકારવું} mean? (\textquotedblleft{}To steer, to drive\textquotedblright{}.) Say it once more.
\end{tcolorbox}
\section[\texorpdfstring{\gu{યાત્રા} (yātrā)}{યાત્રા (yātrā)}]{\gu{યાત્રા} (yātrā) — pilgrimage, journey}

\label{lesson:GU-C211-yatra}



\begin{quote}
Before the new one: say the Gujarati for to ride, then the Gujarati for to steer.
\end{quote}

\subsection*{You'll want to know: \gu{યાત્રા}}
\textbf{\gu{યાત્રા}} — \emph{yātrā} — \textquotedblleft{}pilgrimage, journey\textquotedblright{}.

\begin{quote}
\textbf{\gu{દાદી} \gu{ગયા} \gu{વરસે} \gu{યાત્રાએ} \gu{ગઈ}.} — \emph{dādī gayā varse yātrāe gaī} — \textquotedblleft{}Grandmother went on a pilgrimage last year.\textquotedblright{}
\end{quote}

\begin{cousinweb}
\textbf{\gu{યાત્રા}} (\emph{yātrā}) is Sanskrit \emph{yātrā}, \textquotedblleft{}going, journey, procession\textquotedblright{}, taken over whole as a learned loan, a tatsama. CDIAL has no entry for it.
\end{cousinweb}

\begin{grammarlens}[title={What you've built}]
One more word for everyday life at A2.
\end{grammarlens}

\subsection*{Your turn}
\begin{itemize}
\raggedright
  \item \emph{Say it:} \emph{yātrā}
  \item \emph{Say it:} \emph{yātrā}, once more
  \item \emph{Say it:} \emph{dādī gayā varse yātrāe gaī} — Grandmother went on a pilgrimage last year.
  \item \emph{Recall:} say the Gujarati for to ride, then the Gujarati for to steer, then say \emph{yātrā} again
\end{itemize}

\begin{tcolorbox}[breakable,colback=teal!4,colframe=teal!35!black,title={Before you move on}]
What does \gu{યાત્રા} mean? (\textquotedblleft{}Pilgrimage, journey\textquotedblright{}.) Say it once more.
\end{tcolorbox}
\section[\texorpdfstring{\gu{વાહન} (vāhan)}{વાહન (vāhan)}]{\gu{વાહન} (vāhan) — vehicle}

\label{lesson:GU-C211-vahan}



\begin{quote}
Before the new one: say the Gujarati for to steer, then the Gujarati for pilgrimage.
\end{quote}

\subsection*{You'll want to know: \gu{વાહન}}
\textbf{\gu{વાહન}} — \emph{vāhan} — \textquotedblleft{}vehicle\textquotedblright{}.

\begin{quote}
\textbf{\gu{આ} \gu{વાહન} \gu{જૂનું} \gu{હતું}.} — \emph{ā vāhan jūnũ hatũ} — \textquotedblleft{}This vehicle was old.\textquotedblright{}
\end{quote}

Its inherited cousin is \textbf{\gu{વહાણ}}, \emph{vahāṇ}, 'ship'.

\begin{cousinweb}
\textbf{\gu{વાહન}} (\emph{vāhan}) is Sanskrit \emph{v\'{\={a}}hana}, \textquotedblleft{}any vehicle or draught animal\textquotedblright{}, taken over whole, a tatsama. CDIAL lists its inherited form \emph{vahāṇ}, \textquotedblleft{}ship\textquotedblright{} (CDIAL 11610).
\end{cousinweb}

\begin{grammarlens}[title={What you've built}]
One more word for everyday life at A2.
\end{grammarlens}

\subsection*{Your turn}
\begin{itemize}
\raggedright
  \item \emph{Say it:} \emph{vāhan}
  \item \emph{Say it:} \emph{vāhan}, once more
  \item \emph{Say it:} \emph{ā vāhan jūnũ hatũ} — This vehicle was old.
  \item \emph{Recall:} say the Gujarati for to steer, then the Gujarati for pilgrimage, then say \emph{vāhan} again
\end{itemize}

\begin{tcolorbox}[breakable,colback=teal!4,colframe=teal!35!black,title={Before you move on}]
What does \gu{વાહન} mean? (\textquotedblleft{}Vehicle\textquotedblright{}.) Say it once more.
\end{tcolorbox}
